% !TEX root = ../main.tex
% Scoped primary-source checks and remaining verification: notes/DRAFT10_REVISION.md.
\section{Related Work}
\label{sec:related}

\paragraph{VLA policies and geometric action refinement.}
Flow-based VLAs generate action chunks conditioned on visual and language features~\cite{black2024pi0,intelligence2025pi05,nvidia2025gr00t}.
Their outputs can serve as references for geometric optimization. MoManipVLA transfers VLA-generated end-effector
trajectories to mobile manipulation through optimization with scene geometry~\cite{wu2025momanipvla}. More generally,
cuRobo and cuRoboV2 provide GPU-accelerated motion generation, including collision queries and depth-fused distance
fields~\cite{sundaralingam2023curobo,sundaralingam2026curobov2}. These works motivate our use of a frozen policy and a
geometric refinement layer. Our focus is the information supplied to that layer: how a visually identified target
changes the collision constraints of particular robot links.

\paragraph{Grounding from internal visual representations.}
Selected attention heads in frozen vision-language models can localize referred objects~\cite{kang2025fewheads}.
KNOWS brings this idea to manipulation by using VLA attention to rank objects produced by a fine-tuned YOLOE
segmenter, then applying an ellipsoid-based safety filter~\cite{park2026knows}. It is the closest precedent for using
internal attention to select a contact target. Our target geometry is instead constructed by lifting attention into
depth and growing regions in the observed 3D cloud. This avoids requiring a segmentation model to enumerate candidate
objects and supplies a geometric region for link-specific distance queries. Attention is used as a selection cue;
the quality of the resulting extent must be assessed separately from target identity.

\paragraph{Task-dependent collision geometry.}
VLSA uses a vision-language model and visual grounding to identify obstacles for action
filtering~\cite{hu2026vlsa}. Lan-o3dp combines language-guided object decomposition and segmentation with
collision-aware diffusion-policy execution~\cite{li2025lano3dp}. MoManipVLA's real-world pipeline excludes the target
mask from ESDF construction~\cite{wu2025momanipvla}. Stage-dependent geometry is also used in AutoMoMa, which generates
whole-body trajectories in known scenes and incorporates manipulated objects into the robot's kinematic
representation~\cite{niu2026automoma}. WBAG explicitly models whole-body and grasp-dependent attached geometry for
VLA filtering~\cite{zhen2026wbag}. We build on the need for such contact-dependent representations. Our distinction is
to construct the contact target from VLA attention and observed depth, retain separate distance queries for contact
and non-contact links, and manage these queries during chunk refinement. The current robot query model covers
gripper links and attached-object approximations; it does not provide whole-body coverage.

\paragraph{Reading task state from frozen policies.}
$\pi_{0.5}$ and LoHoVLA incorporate sub-task prediction into their learned policy
structure~\cite{intelligence2025pi05,yang2025lohovla}. Probing studies also recover task progress and success-related
signals from frozen VLA representations~\cite{bhardwaj2026progress,zhang2026frozenvlas}. We use a supervised linear probe
for three manipulation phases and apply its output to target adoption and exception termination. The probe supports
the contact representation; it neither modifies the VLA weights nor generates the policy's action sequence.
Its value beyond a grasp-state rule is an empirical question addressed by a matched-controller comparison.
